\subsection*{4. The 4 Price Thresholds, and Why None Is Universal}

Distinguish 4 operational boundaries:
\begin{enumerate}
  \item A \textbf{short-run incremental-cost floor}, below which each additional transaction may worsen cash contribution.
  \item A \textbf{long-run sustainable floor}, which also pays for capacity, acquisition, management, and replenishment.
  \item A \textbf{demand frontier} at which enough intended buyers actually purchase.
  \item A \textbf{strategic sustainability boundary}, beyond which regret, churn, overpromising, or reputational consequences destroy future value.
\end{enumerate}

The viable commercial zone is the intersection, if 1 exists. Founders sometimes discover that no overlap exists: the price customers accept cannot finance the experience that would make them satisfied. No amount of website polish solves that structural problem. 1 must improve operations, change the offer, choose another segment, alter the buyer/payer relationship, or stop.

A high experimental price can be justified when capacity is scarce and buyer value is strong. But a high price is not evidence of premium quality and a low price is not evidence of weak ambition. Every amount must be evaluated relative to the exchange and the constraint being optimized.

\subsection*{5. 3 Classical Methods, and Where Each Fails}

\textbf{Cost-plus pricing} starts with cost and adds a target markup. It is practical for standardized contracts, regulated settings, and industries where detailed costs are observable. Its blind spot is the buyer's benefit; an expensive-to-produce feature is not automatically valuable.

\textbf{Competition-based pricing} anchors to relevant rivals, which can help meet buyer expectations in comparable categories. Its blind spot is that competitors may have different economics, incentives, and scopes.

\textbf{Value-based pricing} reasons from incremental benefit over alternatives and attempts to capture a sustainable fraction. Its blind spot is information: value may be hard to observe, quantify, attribute, or sell through procurement.

Andreas Hinterhuber's pricing research describes these approaches and the organizational challenges of value-based pricing.\footnote{Andreas Hinterhuber (2008), ``Customer Value-Based Pricing Strategies: Why Companies Resist,'' \emph{Journal of Business Strategy}, 29(4), 41 to 50.} The real founder should use all 3 as evidence, not adopt any as an exclusive dogma. Cost tells you what must be financed; competition tells you what comparisons buyers encounter; value tells you where superior willingness to pay might exist.

An elegant price may be an operational failure when the sales team cannot explain it, the customer's bill cannot be predicted, or the product cannot produce the value being claimed.

\subsection*{6. Economic Value to Customer: A Worked Enterprise Example}

Imagine a customer currently pays ₹2,00,000 annually for a system and spends ₹3,00,000 in associated labor and error handling: ₹5,00,000 total. A proposed alternative charges ₹2,50,000, reduces labor and error cost to ₹1,25,000, and requires ₹35,000 of 1st-year onboarding. The new total is ₹4,10,000. The quoted subscription is more expensive, yet the modeled 1st-year total cost is ₹90,000 lower.

The vendor might legitimately price above the incumbent subscription if the buyer believes the labor and error savings are real. But the model must ask whether reduced minutes actually produce cash savings or useful capacity; whether the baseline contains hidden subsidies; whether changes impose 1-time disruption; and whether a cheaper vendor can produce the same advantage. Those assumptions define the price argument.

This example also demonstrates why the seller should present an auditable \textbf{economic value model}, not a grand ROI claim. Show assumptions, scenarios, and a measurable pilot. When the value is emotional, aesthetic, or social, forced monetary estimates can be misleading; use actual comparative purchasing choices instead.

\subsection*{7. Which Business Objective Does Price Serve?}

A price can maximize short-run profit, cash collections, distribution, adoption, utilization, evidence generation, customer access, contribution per scarce expert-hour, or long-term enterprise value. It cannot generally maximize all simultaneously. A tuition price that fills a class may overwhelm reviewers. An enterprise pilot priced to learn may have lower current margin but help establish a repeatable service. A software penetration price may encourage network formation but create dependence on future cash that never arrives.

Write the objective before testing numbers. If the firm has a hard capacity constraint, contribution per constrained resource may matter more than conversion percentage. If the firm is cash constrained, collection timing may dominate an accounting profit projection. If the firm is trying to build a trusted institution, customer regret and outcome quality are hard strategic constraints, not sentimental metrics appended after growth.
